%==============================================================================
\section{Metodologia}
\label{sec:methodology}

% --- corresponde a 02_Methodology.md; o arco de cinco estágios ---
% O desenvolvimento completo está em 02_Methodology.md; esta seção é a
% condensação em formato de artigo, organizada como adquirir -> transformar
% -> quantificar -> preservar -> inferir.

\subsection{Estágio 1 --- Adquirir}
Reaproveitar os registros sincronizados $(\slk,\Tdie,\Vcore,\tm)$
(\cref{sec:dataset}); excluir o transiente térmico inicial de 1.0~h; manter a
resolução bruta (a discretização é postergada para a etapa de estimação).

\subsection{Estágio 2 --- Transformar}
O acoplamento ambiental \emph{reversível} --- a grandeza que a correção PVT
avalia (\cref{sec:res-correction}) --- é medido sobre sinais com a tendência
retirada (\emph{detrend}, mediana em janela longa), de modo que
$\MI{\scorr}{\Tdie,\Vcore}$ não confunda a deriva lenta de envelhecimento com
acoplamento ambiental. O número-manchete $\MI{\slk}{\Tdie,\Vcore}$ da
\cref{sec:res-mi}, por construção, \textbf{não} é detrendizado: ele é
calculado sobre a mesma série bruta pós-\emph{soak} usada para o $R^2$ do
trabalho anterior, precisamente para permitir a comparação
maçã-com-maçã da \cref{sec:res-mi}. Padronizar $\Tdie,\Vcore$. Tratar a
correlação serial via (i)~o tempo de autocorrelação integrado $\tauint$ e um
tamanho efetivo de amostra $\neff=n/(2\tauint)$; (ii)~intervalos de confiança
por \emph{moving-block bootstrap}; (iii)~decimação como verificação cruzada.
Duas famílias de estimadores são comparadas: \emph{plug-in} com
\emph{binning} e correção de viés de Miller--Madow, e $k$-vizinhos-mais-próximos
(Kozachenko--Leonenko / KSG); as duas concordam na conclusão qualitativa
(acoplamento muito acima do piso nulo; excesso positivo sobre o
equivalente Gaussiano), mas divergem além do intervalo de confiança do
estimador com \emph{binning} --- essa diferença sistemática entre
estimadores é caracterizada, não escondida, na \cref{sec:res-mi}.

\subsection{Estágio 3 --- Quantificar}
Estimar $\Ent{\slk}$, $\Ent{\slk\mid\Tdie,\Vcore}$, as informações mútuas por
canal $\MI{\slk}{\Tdie}$, $\MI{\slk}{\Vcore}$, a conjunta
$\MI{\slk}{\Tdie,\Vcore}$ e a fração de informação; comparar regimes;
calcular $\KL{\cdot}{\cdot}$ entre bancadas (\cref{eq:mi}).

\subsection{Estágio 4 --- Preservar (correção PVT como recuperação de informação)}
Ajustar $\hat{\dpvt}(\Tdie,\Vcore)$ (OLS linear; GLS/MLE não-linear;
Bayesiano) e formar $\scorr=\slk-\hat{\dpvt}$. Avaliar cada correção pela MI
residual \emph{reversível} (detrendizada) $\MI{\scorr}{\Tdie,\Vcore}\to 0$ e
pelo desvio-padrão residual em relação ao piso de quantização; selecionar o
modelo por MDL. Classificar $\Tdie$ vs.\ $\Vcore$ por contribuição de
informação (seleção de atributos). Como um dual exploratório, executar
separação cega por ICA/negentropia e avaliá-la pela MI residual entre
componentes.

\subsection{Estágio 5 --- Inferir (capacidade e RUL)}
Converter o ruído residual pós-correção em estados de envelhecimento
resolvíveis via \cref{eq:capacity}; modelar o alarme por limiar como um BSC e
reportar $\Cbsc=1-\Hb(p)$ (\cref{eq:bsc}). Derivar o limite de
detectabilidade de Cram\'er--Rao sobre a inclinação de envelhecimento usando
$\neff$, complementado por uma posterior Bayesiana. Como uma segunda rota de
estimação sequencial para a mesma questão (\cref{sec:th-wiener}), fazer a
média em blocos de $\scorr$ em $\sim\!\neff$ observações
quase-independentes (comprimento de bloco $2\tauint$, ou seja, uma unidade de
$\neff$ por bloco por construção), ajustar a difusão de Wiener $\sigma_B^2$ e
o ruído de observação por bloco $R$ pelo método dos momentos (variância dos
incrementos de média de bloco vs.\ defasagem), ajustar o ruído de taxa
$\sigma_{\mathrm{rate}}^2$ por MLE unidimensional sobre a verossimilhança de
inovações do filtro de Kalman, e executar um suavizador Kalman/RTS para
obter uma trajetória de envelhecimento denoised $\hat\Xlat(t)$ com bandas de
credibilidade e a distribuição de RUL Gaussiana-inversa da
\cref{eq:invgauss} em dois limiares já calculados
($D=\sigma_{\text{noise}}$, $D=\sigma_{\text{signal}}$). Por fim, enquadrar a
``melhor condição de \emph{burn-in}'' como a minimização de
$\MI{\slk}{\Tdie,\Vcore}$ sujeita a entregar o fator de aceleração de
Arrhenius exigido.

\subsection{Validação}
Comparação entre estimadores independentes (Miller--Madow com
\emph{binning} vs.\ KSG), com a diferença entre eles reportada como faixa de
incerteza de estimador, não suprimida; ordenação consistente entre bancadas
com o ranqueamento de $\sigma$ publicado; testes surrogate/nulos por
randomização de fase para calibrar o piso de viés de MI; \emph{dither}
ligado/desligado para controlar artefatos de quantização; sensibilidade ao
número de \emph{bins} de $\Tdie,\Vcore$ reportada explicitamente
(\cref{sec:res-mi}).
